% GENERATED FILE. Edit canonical lessons, then run npm run generate:books.
% canonical-source-hash: fnv1a64:a23deca3787ed718
% canonical-lessons: FR-C530-vaisselle, FR-C530-evier, FR-C530-rincer, FR-C530-essuyer, FR-C530-torchon

\chapter{Faire la vaisselle}
\label{ch:fr-faire-la-vaisselle}
\hlchaptermodality{530}

\begin{chapteropening}
\textbf{By the end of this chapter:} \emph{I can talk about doing the dishes.}

Complete the last lesson of chapter 530: I can talk about doing the dishes.
\end{chapteropening}

\section[la vaisselle]{la vaisselle — dishes, crockery}

\label{lesson:FR-C530-vaisselle}



\begin{quote}
Before the new one: say the French for ladle, then the French for lid.
\end{quote}

\subsection*{You'll want to know: la vaisselle}
\textbf{la vaisselle} — \textquotedblleft{}dishes, crockery\textquotedblright{}.

\begin{quote}
Ce soir, c'est mon frère qui fait la vaisselle.
\end{quote}

\begin{quote}
\emph{Tonight it is my brother who does the dishes.}
\end{quote}

\emph{Faire la vaisselle} — to do the washing-up. In Switzerland also \emph{relaver}.

\begin{cousinweb}
\emph{la vaisselle} is inherited, through Old French \emph{vaissele}, from Vulgar Latin \emph{vascella} (EtymDB), from Latin \emph{vāsculum} \textquotedblleft{}a small vessel\textquotedblright{}, built on \emph{vās} \textquotedblleft{}a vessel\textquotedblright{} (Lewis \& Short).
\end{cousinweb}

\begin{grammarlens}[title={What you've built}]
One more word for telling, explaining and describing at B1.
\end{grammarlens}

\subsection*{Your turn}
\begin{itemize}
\raggedright
  \item \emph{Say it:} \emph{la vaisselle}
  \item \emph{Say it:} \emph{la vaisselle}, once more
  \item \emph{Say it:} a sentence of your own with \emph{la vaisselle}
  \item \emph{Recall:} say the French for ladle, then the French for lid, then say \emph{la vaisselle} again
\end{itemize}

\begin{tcolorbox}[breakable,colback=teal!4,colframe=teal!35!black,title={Before you move on}]
What does la vaisselle mean? (\textquotedblleft{}Dishes, crockery\textquotedblright{}.) Say it once more.
\end{tcolorbox}
\section[l'évier]{l'évier — (kitchen) sink}

\label{lesson:FR-C530-evier}



\begin{quote}
Before the new one: say the French for lid, then the French for dishes.
\end{quote}

\subsection*{You'll want to know: l'évier}
\textbf{l'évier} — \textquotedblleft{}(kitchen) sink\textquotedblright{}.

\begin{quote}
Les assiettes sales sont dans l'évier.
\end{quote}

\begin{quote}
\emph{The dirty plates are in the sink.}
\end{quote}

Masculine. A bathroom basin is \emph{le lavabo}.

\begin{cousinweb}
\emph{l'évier} is inherited from Latin \emph{aquārium} \textquotedblleft{}a place for water\textquotedblright{} (EtymDB), built on \emph{aqua} \textquotedblleft{}water\textquotedblright{} (Lewis \& Short).
\end{cousinweb}

\begin{grammarlens}[title={What you've built}]
One more word for telling, explaining and describing at B1.
\end{grammarlens}

\subsection*{Your turn}
\begin{itemize}
\raggedright
  \item \emph{Say it:} \emph{l'évier}
  \item \emph{Say it:} \emph{l'évier}, once more
  \item \emph{Say it:} a sentence of your own with \emph{l'évier}
  \item \emph{Recall:} say the French for lid, then the French for dishes, then say \emph{l'évier} again
\end{itemize}

\begin{tcolorbox}[breakable,colback=teal!4,colframe=teal!35!black,title={Before you move on}]
What does l'évier mean? (\textquotedblleft{}(kitchen) sink\textquotedblright{}.) Say it once more.
\end{tcolorbox}
\section[rincer]{rincer — to rinse}

\label{lesson:FR-C530-rincer}



\begin{quote}
Before the new one: say the French for dishes, then the French for (kitchen) sink.
\end{quote}

\subsection*{You'll want to know: rincer}
\textbf{rincer} — \textquotedblleft{}to rinse\textquotedblright{}.

\begin{quote}
Je rince les verres à l'eau chaude.
\end{quote}

\begin{quote}
\emph{I rinse the glasses in hot water.}
\end{quote}

Spelling: \emph{nous rinçons}.

\begin{cousinweb}
\emph{rincer} goes back to Old French \emph{reïncier}. The origin is not settled: a Vulgar Latin \emph{recentiāre} \textquotedblleft{}to refresh\textquotedblright{}, from \emph{recēns} \textquotedblleft{}fresh\textquotedblright{} (Lewis \& Short), is often proposed, and EtymDB compares Old Norse \emph{hreinsa} \textquotedblleft{}to cleanse\textquotedblright{}.
\end{cousinweb}

\begin{grammarlens}[title={What you've built}]
One more word for telling, explaining and describing at B1.
\end{grammarlens}

\subsection*{Your turn}
\begin{itemize}
\raggedright
  \item \emph{Say it:} \emph{rincer}
  \item \emph{Say it:} \emph{rincer}, once more
  \item \emph{Say it:} a sentence of your own with \emph{rincer}
  \item \emph{Recall:} say the French for dishes, then the French for (kitchen) sink, then say \emph{rincer} again
\end{itemize}

\begin{tcolorbox}[breakable,colback=teal!4,colframe=teal!35!black,title={Before you move on}]
What does rincer mean? (\textquotedblleft{}To rinse\textquotedblright{}.) Say it once more.
\end{tcolorbox}
\section[essuyer]{essuyer — to wipe, to dry}

\label{lesson:FR-C530-essuyer}



\begin{quote}
Before the new one: say the French for (kitchen) sink, then the French for to rinse.
\end{quote}

\subsection*{You'll want to know: essuyer}
\textbf{essuyer} — \textquotedblleft{}to wipe, to dry\textquotedblright{}.

\begin{quote}
Il essuie les assiettes après le repas.
\end{quote}

\begin{quote}
\emph{He dries the plates after the meal.}
\end{quote}

\emph{J'essuie}, \emph{nous essuyons}. \emph{S'essuyer les mains} — to dry one's hands.

\begin{cousinweb}
\emph{essuyer} is inherited from Latin \emph{exsūcāre} \textquotedblleft{}to squeeze the juice out\textquotedblright{}, built on \emph{sūcus} \textquotedblleft{}juice\textquotedblright{} (Lewis \& Short).
\end{cousinweb}

\begin{grammarlens}[title={What you've built}]
One more word for telling, explaining and describing at B1.
\end{grammarlens}

\subsection*{Your turn}
\begin{itemize}
\raggedright
  \item \emph{Say it:} \emph{essuyer}
  \item \emph{Say it:} \emph{essuyer}, once more
  \item \emph{Say it:} a sentence of your own with \emph{essuyer}
  \item \emph{Recall:} say the French for (kitchen) sink, then the French for to rinse, then say \emph{essuyer} again
\end{itemize}

\begin{tcolorbox}[breakable,colback=teal!4,colframe=teal!35!black,title={Before you move on}]
What does essuyer mean? (\textquotedblleft{}To wipe, to dry\textquotedblright{}.) Say it once more.
\end{tcolorbox}
\section[le torchon]{le torchon — tea towel, dishcloth}

\label{lesson:FR-C530-torchon}



\begin{quote}
Before the new one: say the French for to rinse, then the French for to wipe.
\end{quote}

\subsection*{You'll want to know: le torchon}
\textbf{le torchon} — \textquotedblleft{}tea towel, dishcloth\textquotedblright{}.

\begin{quote}
Le torchon est accroché près de l'évier.
\end{quote}

\begin{quote}
\emph{The tea towel is hanging by the sink.}
\end{quote}

In Belgium \emph{l'essuie}. \emph{Le torchon brûle} — there is trouble between them.

\begin{cousinweb}
\emph{le torchon} is built on \emph{torche} \textquotedblleft{}a twist of straw\textquotedblright{}, which goes back to Vulgar Latin \emph{torca}, from Latin \emph{torquēre} \textquotedblleft{}to twist\textquotedblright{} (Lewis \& Short).
\end{cousinweb}

\begin{grammarlens}[title={What you've built}]
One more word for telling, explaining and describing at B1.
\end{grammarlens}

\subsection*{Your turn}
\begin{itemize}
\raggedright
  \item \emph{Say it:} \emph{le torchon}
  \item \emph{Say it:} \emph{le torchon}, once more
  \item \emph{Say it:} a sentence of your own with \emph{le torchon}
  \item \emph{Recall:} say the French for to rinse, then the French for to wipe, then say \emph{le torchon} again
\end{itemize}

\begin{tcolorbox}[breakable,colback=teal!4,colframe=teal!35!black,title={Before you move on}]
What does le torchon mean? (\textquotedblleft{}Tea towel, dishcloth\textquotedblright{}.) Say it once more.
\end{tcolorbox}
